\documentclass[aspectratio=169,11pt]{beamer}
\usetheme{Madrid}
\usecolortheme{default}
\setbeamertemplate{navigation symbols}{}
\setbeamertemplate{footline}[frame number]
\usepackage{enumitem}
\usepackage{amsmath,amssymb}
\usepackage{booktabs}
\usepackage{array}

\title{Databricks ML Associate}
\subtitle{Questions 401--500 (Fifth 100)}
\author{Practice Exam Deck}
\date{}

\begin{document}
	
	\begin{frame}
		\titlepage
	\end{frame}
	
	\begin{frame}{How to Use This Deck}
		\begin{itemize}
			\item 100 new questions (Q401--Q500), none repeated from Q1--Q400.
			\item Each question slide has four options (A--D).
			\item The next slide gives the correct answer and explains why the right option is right and why each wrong option is wrong.
			\item This set focuses on MLflow search syntax, model signatures, pipeline persistence, Hyperopt internals, Feature Store batch/stream trade-offs, Unity Catalog lineage, Asset Bundles, MLflow evaluate, monitoring, and distributed training.
		\end{itemize}
	\end{frame}
	
	% ---------------------------------------------------------------
	\begin{frame}{Question 401}
		Which MLflow search filter string finds runs where the metric \texttt{rmse} is less than 1?
		\begin{enumerate}[label=\Alph*.]
			\item \texttt{"metrics.rmse < 1"}
			\item \texttt{"rmse < 1"}
			\item \texttt{"metric\_rmse < 1"}
			\item \texttt{"metrics.rmse =< 1"}
		\end{enumerate}
	\end{frame}
	\begin{frame}{Answer 401}
		\textbf{Correct: A} --- filters use \texttt{metrics.<name>}.\\
		\vspace{0.5em}
		\textbf{Wrong:}
		\begin{itemize}
			\item B omits the \texttt{metrics.} prefix.
			\item C is not the syntax.
			\item D uses an invalid operator.
		\end{itemize}
	\end{frame}
	
	% ---------------------------------------------------------------
	\begin{frame}{Question 402}
		Which MLflow search filter matches runs whose parameter \texttt{max\_depth} equals 5?
		\begin{enumerate}[label=\Alph*.]
			\item \texttt{"params.max\_depth = '5'"}
			\item \texttt{"param.max\_depth = 5"}
			\item \texttt{"params.max\_depth == 5"}
			\item \texttt{"max\_depth = 5"}
		\end{enumerate}
	\end{frame}
	\begin{frame}{Answer 402}
		\textbf{Correct: A} --- parameter filters use \texttt{params.<name>} and quoted values.\\
		\vspace{0.5em}
		\textbf{Wrong:}
		\begin{itemize}
			\item B uses the wrong prefix.
			\item C uses \texttt{==} and no quotes.
			\item D omits the prefix.
		\end{itemize}
	\end{frame}
	
	% ---------------------------------------------------------------
	\begin{frame}{Question 403}
		Which MLflow API marks a run as finished?
		\begin{enumerate}[label=\Alph*.]
			\item \texttt{mlflow.end\_run()}
			\item \texttt{mlflow.finish\_run()}
			\item \texttt{mlflow.close\_run()}
			\item \texttt{mlflow.stop\_run()}
		\end{enumerate}
	\end{frame}
	\begin{frame}{Answer 403}
		\textbf{Correct: A} --- \texttt{end\_run()} finishes a run.\\
		\vspace{0.5em}
		\textbf{Wrong:}
		\begin{itemize}
			\item B, C, D are not the API.
		\end{itemize}
	\end{frame}
	
	% ---------------------------------------------------------------
	\begin{frame}{Question 404}
		Which MLflow API deletes a run?
		\begin{enumerate}[label=\Alph*.]
			\item \texttt{mlflow.delete\_run()}
			\item \texttt{mlflow.remove\_run()}
			\item \texttt{mlflow.drop\_run()}
			\item \texttt{mlflow.end\_run(delete=True)}
		\end{enumerate}
	\end{frame}
	\begin{frame}{Answer 404}
		\textbf{Correct: A} --- \texttt{delete\_run()} deletes a run.\\
		\vspace{0.5em}
		\textbf{Wrong:}
		\begin{itemize}
			\item B, C, D are not the API.
		\end{itemize}
	\end{frame}
	
	% ---------------------------------------------------------------
	\begin{frame}{Question 405}
		What does \texttt{mlflow.active\_run()} return?
		\begin{enumerate}[label=\Alph*.]
			\item The currently active run object
			\item A list of all runs
			\item A metric value
			\item A model
		\end{enumerate}
	\end{frame}
	\begin{frame}{Answer 405}
		\textbf{Correct: A} --- it returns the active run.\\
		\vspace{0.5em}
		\textbf{Wrong:}
		\begin{itemize}
			\item B, C, D are not returned.
		\end{itemize}
	\end{frame}
	
	% ---------------------------------------------------------------
	\begin{frame}{Question 406}
		What is the purpose of \texttt{mlflow.set\_tracking\_uri()}?
		\begin{enumerate}[label=\Alph*.]
			\item It sets the tracking backend location
			\item It sets the model registry name
			\item It sets the cluster size
			\item It sets the feature table
		\end{enumerate}
	\end{frame}
	\begin{frame}{Answer 406}
		\textbf{Correct: A} --- it points MLflow to a tracking backend.\\
		\vspace{0.5em}
		\textbf{Wrong:}
		\begin{itemize}
			\item B, C, D are unrelated.
		\end{itemize}
	\end{frame}
	
	% ---------------------------------------------------------------
	\begin{frame}{Question 407}
		What does \texttt{mlflow.models.infer\_signature()} do?
		\begin{enumerate}[label=\Alph*.]
			\item It infers input/output schema from sample data
			\item It trains the model
			\item It registers the model
			\item It deploys the model
		\end{enumerate}
	\end{frame}
	\begin{frame}{Answer 407}
		\textbf{Correct: A} --- it infers a model signature.\\
		\vspace{0.5em}
		\textbf{Wrong:}
		\begin{itemize}
			\item B, C, D are unrelated.
		\end{itemize}
	\end{frame}
	
	% ---------------------------------------------------------------
	\begin{frame}{Question 408}
		Why is a model signature useful for serving?
		\begin{enumerate}[label=\Alph*.]
			\item It validates request and response schemas
			\item It trains the model
			\item It stores features
			\item It deletes versions
		\end{enumerate}
	\end{frame}
	\begin{frame}{Answer 408}
		\textbf{Correct: A} --- signatures validate schemas.\\
		\vspace{0.5em}
		\textbf{Wrong:}
		\begin{itemize}
			\item B, C, D are false.
		\end{itemize}
	\end{frame}
	
	% ---------------------------------------------------------------
	\begin{frame}{Question 409}
		What does \texttt{mlflow.models.set\_signature()} do?
		\begin{enumerate}[label=\Alph*.]
			\item It attaches a signature to a logged model
			\item It creates a signature from scratch
			\item It deletes a signature
			\item It registers a signature
		\end{enumerate}
	\end{frame}
	\begin{frame}{Answer 409}
		\textbf{Correct: A} --- it sets a model's signature.\\
		\vspace{0.5em}
		\textbf{Wrong:}
		\begin{itemize}
			\item B, C, D are not its behavior.
		\end{itemize}
	\end{frame}
	
	% ---------------------------------------------------------------
	\begin{frame}{Question 410}
		What does \texttt{PipelineModel.save()} do in Spark ML?
		\begin{enumerate}[label=\Alph*.]
			\item Persists the fitted pipeline to a path
			\item Trains the pipeline
			\item Registers the model
			\item Serves the model
		\end{enumerate}
	\end{frame}
	\begin{frame}{Answer 410}
		\textbf{Correct: A} --- \texttt{save()} persists a fitted pipeline.\\
		\vspace{0.5em}
		\textbf{Wrong:}
		\begin{itemize}
			\item B, C, D are other operations.
		\end{itemize}
	\end{frame}
	
	% ---------------------------------------------------------------
	\begin{frame}{Question 411}
		What does \texttt{PipelineModel.load()} do in Spark ML?
		\begin{enumerate}[label=\Alph*.]
			\item Loads a saved pipeline model from a path
			\item Trains a pipeline
			\item Registers a model
			\item Serves a model
		\end{enumerate}
	\end{frame}
	\begin{frame}{Answer 411}
		\textbf{Correct: A} --- \texttt{load()} restores a saved model.\\
		\vspace{0.5em}
		\textbf{Wrong:}
		\begin{itemize}
			\item B, C, D are other operations.
		\end{itemize}
	\end{frame}
	
	% ---------------------------------------------------------------
	\begin{frame}{Question 412}
		Why save the entire pipeline instead of only the estimator?
		\begin{enumerate}[label=\Alph*.]
			\item It preserves preprocessing steps needed at inference
			\item It trains faster
			\item It reduces storage
			\item It registers the model
		\end{enumerate}
	\end{frame}
	\begin{frame}{Answer 412}
		\textbf{Correct: A} --- the pipeline includes preprocessing.\\
		\vspace{0.5em}
		\textbf{Wrong:}
		\begin{itemize}
			\item B, C, D are false.
		\end{itemize}
	\end{frame}
	
	% ---------------------------------------------------------------
	\begin{frame}{Question 413}
		What does \texttt{Pipeline.fit()} return?
		\begin{enumerate}[label=\Alph*.]
			\item A \texttt{PipelineModel} with fitted stages
			\item A DataFrame
			\item A metric
			\item A cluster
		\end{enumerate}
	\end{frame}
	\begin{frame}{Answer 413}
		\textbf{Correct: A} --- \texttt{fit()} returns a \texttt{PipelineModel}.\\
		\vspace{0.5em}
		\textbf{Wrong:}
		\begin{itemize}
			\item B, C, D are not returned.
		\end{itemize}
	\end{frame}
	
	% ---------------------------------------------------------------
	\begin{frame}{Question 414}
		What happens if a pipeline stage is a transformer only?
		\begin{enumerate}[label=\Alph*.]
			\item It is applied without fitting
			\item It is trained
			\item It is skipped
			\item It raises an error
		\end{enumerate}
	\end{frame}
	\begin{frame}{Answer 414}
		\textbf{Correct: A} --- transformers are applied, not fit.\\
		\vspace{0.5em}
		\textbf{Wrong:}
		\begin{itemize}
			\item B is for estimators.
			\item C, D are false.
		\end{itemize}
	\end{frame}
	
	% ---------------------------------------------------------------
	\begin{frame}{Question 415}
		What is an estimator in Spark ML?
		\begin{enumerate}[label=\Alph*.]
			\item An algorithm that fits data and produces a model
			\item A fitted model
			\item A DataFrame
			\item A metric
		\end{enumerate}
	\end{frame}
	\begin{frame}{Answer 415}
		\textbf{Correct: A} --- estimators produce models via \texttt{fit()}.\\
		\vspace{0.5em}
		\textbf{Wrong:}
		\begin{itemize}
			\item B is a transformer/model.
			\item C, D are not estimators.
		\end{itemize}
	\end{frame}
	
	% ---------------------------------------------------------------
	\begin{frame}{Question 416}
		What is a transformer in Spark ML?
		\begin{enumerate}[label=\Alph*.]
			\item An object that transforms a DataFrame into another DataFrame
			\item An algorithm that produces a model
			\item A metric
			\item A cluster
		\end{enumerate}
	\end{frame}
	\begin{frame}{Answer 416}
		\textbf{Correct: A} --- transformers map DataFrames.\\
		\vspace{0.5em}
		\textbf{Wrong:}
		\begin{itemize}
			\item B is an estimator.
			\item C, D are unrelated.
		\end{itemize}
	\end{frame}
	
	% ---------------------------------------------------------------
	\begin{frame}{Question 417}
		What does \texttt{hp.quniform()} define in Hyperopt?
		\begin{enumerate}[label=\Alph*.]
			\item A uniform distribution quantized to integer steps
			\item A categorical choice
			\item A log-uniform distribution
			\item A constant
		\end{enumerate}
	\end{frame}
	\begin{frame}{Answer 417}
		\textbf{Correct: A} --- \texttt{hp.quniform} is a quantized uniform space.\\
		\vspace{0.5em}
		\textbf{Wrong:}
		\begin{itemize}
			\item B is \texttt{hp.choice}.
			\item C is \texttt{hp.loguniform}.
			\item D is a constant.
		\end{itemize}
	\end{frame}
	
	% ---------------------------------------------------------------
	\begin{frame}{Question 418}
		What does \texttt{hp.randint()} define in Hyperopt?
		\begin{enumerate}[label=\Alph*.]
			\item A random integer in a range
			\item A categorical choice
			\item A log-uniform space
			\item A constant
		\end{enumerate}
	\end{frame}
	\begin{frame}{Answer 418}
		\textbf{Correct: A} --- \texttt{hp.randint} samples integers.\\
		\vspace{0.5em}
		\textbf{Wrong:}
		\begin{itemize}
			\item B, C, D are other spaces.
		\end{itemize}
	\end{frame}
	
	% ---------------------------------------------------------------
	\begin{frame}{Question 419}
		What does \texttt{hp.loguniform()} define in Hyperopt?
		\begin{enumerate}[label=\Alph*.]
			\item A log-uniform distribution between bounds
			\item A uniform distribution
			\item A categorical choice
			\item A constant
		\end{enumerate}
	\end{frame}
	\begin{frame}{Answer 419}
		\textbf{Correct: A} --- \texttt{hp.loguniform} is log-uniform.\\
		\vspace{0.5em}
		\textbf{Wrong:}
		\begin{itemize}
			\item B is \texttt{hp.uniform}.
			\item C is \texttt{hp.choice}.
			\item D is a constant.
		\end{itemize}
	\end{frame}
	
	% ---------------------------------------------------------------
	\begin{frame}{Question 420}
		What does the objective function return in Hyperopt?
		\begin{enumerate}[label=\Alph*.]
			\item A loss value to minimize and optional metadata
			\item A model
			\item A DataFrame
			\item A cluster
		\end{enumerate}
	\end{frame}
	\begin{frame}{Answer 420}
		\textbf{Correct: A} --- the objective returns a loss (and optional status/dict).\\
		\vspace{0.5em}
		\textbf{Wrong:}
		\begin{itemize}
			\item B, C, D are not returned.
		\end{itemize}
	\end{frame}
	
	% ---------------------------------------------------------------
	\begin{frame}{Question 421}
		What does \texttt{STATUS\_OK} mean in a Hyperopt objective return?
		\begin{enumerate}[label=\Alph*.]
			\item The trial completed successfully
			\item The trial failed
			\item The trial was pruned
			\item The trial is running
		\end{enumerate}
	\end{frame}
	\begin{frame}{Answer 421}
		\textbf{Correct: A} --- \texttt{STATUS\_OK} marks success.\\
		\vspace{0.5em}
		\textbf{Wrong:}
		\begin{itemize}
			\item B is \texttt{STATUS\_FAIL}.
			\item C is pruned.
			\item D is running.
		\end{itemize}
	\end{frame}
	
	% ---------------------------------------------------------------
	\begin{frame}{Question 422}
		What is the benefit of pruning trials in Hyperopt?
		\begin{enumerate}[label=\Alph*.]
			\item It stops poor trials early to save compute
			\item It improves accuracy
			\item It registers models
			\item It serves models
		\end{enumerate}
	\end{frame}
	\begin{frame}{Answer 422}
		\textbf{Correct: A} --- pruning saves compute on poor trials.\\
		\vspace{0.5em}
		\textbf{Wrong:}
		\begin{itemize}
			\item B, C, D are false.
		\end{itemize}
	\end{frame}
	
	% ---------------------------------------------------------------
	\begin{frame}{Question 423}
		What does \texttt{SparkTrials(parallelism=N)} control?
		\begin{enumerate}[label=\Alph*.]
			\item The number of concurrent Hyperopt trials
			\item The number of folds
			\item The number of trees
			\item The number of features
		\end{enumerate}
	\end{frame}
	\begin{frame}{Answer 423}
		\textbf{Correct: A} --- \texttt{parallelism} sets concurrent trials.\\
		\vspace{0.5em}
		\textbf{Wrong:}
		\begin{itemize}
			\item B, C, D are other settings.
		\end{itemize}
	\end{frame}
	
	% ---------------------------------------------------------------
	\begin{frame}{Question 424}
		What does \texttt{fmin} return?
		\begin{enumerate}[label=\Alph*.]
			\item The best hyperparameters and the best loss
			\item A trained model
			\item A DataFrame
			\item A cluster
		\end{enumerate}
	\end{frame}
	\begin{frame}{Answer 424}
		\textbf{Correct: A} --- \texttt{fmin} returns best params and loss.\\
		\vspace{0.5em}
		\textbf{Wrong:}
		\begin{itemize}
			\item B, C, D are not returned.
		\end{itemize}
	\end{frame}
	
	% ---------------------------------------------------------------
	\begin{frame}{Question 425}
		What is the difference between a feature table and a training set?
		\begin{enumerate}[label=\Alph*.]
			\item Feature tables store features; training sets join features with labels
			\item They are identical
			\item Training sets store data; feature tables are views
			\item Feature tables are models
		\end{enumerate}
	\end{frame}
	\begin{frame}{Answer 425}
		\textbf{Correct: A} --- training sets join features and labels.\\
		\vspace{0.5em}
		\textbf{Wrong:}
		\begin{itemize}
			\item B, C, D are false.
		\end{itemize}
	\end{frame}
	
	% ---------------------------------------------------------------
	\begin{frame}{Question 426}
		What is the purpose of an online feature store in real-time serving?
		\begin{enumerate}[label=\Alph*.]
			\item It provides low-latency feature lookups
			\item It trains models
			\item It registers models
			\item It stores notebooks
		\end{enumerate}
	\end{frame}
	\begin{frame}{Answer 426}
		\textbf{Correct: A} --- online stores serve low-latency features.\\
		\vspace{0.5em}
		\textbf{Wrong:}
		\begin{itemize}
			\item B, C, D are unrelated.
		\end{itemize}
	\end{frame}
	
	% ---------------------------------------------------------------
	\begin{frame}{Question 427}
		What is the purpose of an offline feature store?
		\begin{enumerate}[label=\Alph*.]
			\item It stores features for batch training and scoring
			\item It serves real-time requests
			\item It registers models
			\item It stores notebooks
		\end{enumerate}
	\end{frame}
	\begin{frame}{Answer 427}
		\textbf{Correct: A} --- offline stores support batch workflows.\\
		\vspace{0.5em}
		\textbf{Wrong:}
		\begin{itemize}
			\item B is online.
			\item C, D are unrelated.
		\end{itemize}
	\end{frame}
	
	% ---------------------------------------------------------------
	\begin{frame}{Question 428}
		What is a batch feature table used for?
		\begin{enumerate}[label=\Alph*.]
			\item Providing features for offline training and batch scoring
			\item Serving low-latency inference
			\item Registering models
			\item Storing notebooks
		\end{enumerate}
	\end{frame}
	\begin{frame}{Answer 428}
		\textbf{Correct: A} --- batch feature tables support offline use.\\
		\vspace{0.5em}
		\textbf{Wrong:}
		\begin{itemize}
			\item B is online.
			\item C, D are unrelated.
		\end{itemize}
	\end{frame}
	
	% ---------------------------------------------------------------
	\begin{frame}{Question 429}
		What is a streaming feature table used for?
		\begin{enumerate}[label=\Alph*.]
			\item Continuously updated features for streaming inference
			\item Offline batch training
			\item Model registration
			\item Notebook storage
		\end{enumerate}
	\end{frame}
	\begin{frame}{Answer 429}
		\textbf{Correct: A} --- streaming tables support live features.\\
		\vspace{0.5em}
		\textbf{Wrong:}
		\begin{itemize}
			\item B is offline.
			\item C, D are unrelated.
		\end{itemize}
	\end{frame}
	
	% ---------------------------------------------------------------
	\begin{frame}{Question 430}
		Why is a timestamp column important in feature tables?
		\begin{enumerate}[label=\Alph*.]
			\item It enables point-in-time joins
			\item It trains the model
			\item It registers the model
			\item It serves the model
		\end{enumerate}
	\end{frame}
	\begin{frame}{Answer 430}
		\textbf{Correct: A} --- timestamps enable point-in-time joins.\\
		\vspace{0.5em}
		\textbf{Wrong:}
		\begin{itemize}
			\item B, C, D are unrelated.
		\end{itemize}
	\end{frame}
	
	% ---------------------------------------------------------------
	\begin{frame}{Question 431}
		What is Unity Catalog lineage?
		\begin{enumerate}[label=\Alph*.]
			\item A record of data and model dependencies
			\item A model
			\item A feature table
			\item A notebook
		\end{enumerate}
	\end{frame}
	\begin{frame}{Answer 431}
		\textbf{Correct: A} --- lineage tracks dependencies.\\
		\vspace{0.5em}
		\textbf{Wrong:}
		\begin{itemize}
			\item B, C, D are unrelated.
		\end{itemize}
	\end{frame}
	
	% ---------------------------------------------------------------
	\begin{frame}{Question 432}
		Why is lineage useful for ML models?
		\begin{enumerate}[label=\Alph*.]
			\item It shows upstream data and downstream consumers
			\item It trains models
			\item It serves models
			\item It registers models
		\end{enumerate}
	\end{frame}
	\begin{frame}{Answer 432}
		\textbf{Correct: A} --- lineage reveals upstream/downstream dependencies.\\
		\vspace{0.5em}
		\textbf{Wrong:}
		\begin{itemize}
			\item B, C, D are unrelated.
		\end{itemize}
	\end{frame}
	
	% ---------------------------------------------------------------
	\begin{frame}{Question 433}
		What does a Databricks Asset Bundle define?
		\begin{enumerate}[label=\Alph*.]
			\item Infrastructure and jobs as code
			\item A model
			\item A feature table
			\item A notebook
		\end{enumerate}
	\end{frame}
	\begin{frame}{Answer 433}
		\textbf{Correct: A} --- Asset Bundles define infrastructure as code.\\
		\vspace{0.5em}
		\textbf{Wrong:}
		\begin{itemize}
			\item B, C, D are not bundles.
		\end{itemize}
	\end{frame}
	
	% ---------------------------------------------------------------
	\begin{frame}{Question 434}
		Which file format is used to define a Databricks Asset Bundle?
		\begin{enumerate}[label=\Alph*.]
			\item YAML
			\item JSON
			\item XML
			\item CSV
		\end{enumerate}
	\end{frame}
	\begin{frame}{Answer 434}
		\textbf{Correct: A} --- bundles use YAML.\\
		\vspace{0.5em}
		\textbf{Wrong:}
		\begin{itemize}
			\item B, C, D are not the bundle format.
		\end{itemize}
	\end{frame}
	
	% ---------------------------------------------------------------
	\begin{frame}{Question 435}
		What does \texttt{mlflow.evaluate()} do?
		\begin{enumerate}[label=\Alph*.]
			\item It evaluates a model against a dataset and logs metrics
			\item It trains a model
			\item It registers a model
			\item It serves a model
		\end{enumerate}
	\end{frame}
	\begin{frame}{Answer 435}
		\textbf{Correct: A} --- \texttt{evaluate()} computes and logs metrics.\\
		\vspace{0.5em}
		\textbf{Wrong:}
		\begin{itemize}
			\item B, C, D are other operations.
		\end{itemize}
	\end{frame}
	
	% ---------------------------------------------------------------
	\begin{frame}{Question 436}
		Which evaluator type does \texttt{mlflow.evaluate()} accept for classification?
		\begin{enumerate}[label=\Alph*.]
			\item A classifier evaluator with metrics such as F1 and ROC-AUC
			\item A clustering evaluator
			\item A regression evaluator
			\item A summarization evaluator
		\end{enumerate}
	\end{frame}
	\begin{frame}{Answer 436}
		\textbf{Correct: A} --- classifier evaluators compute classification metrics.\\
		\vspace{0.5em}
		\textbf{Wrong:}
		\begin{itemize}
			\item B, C, D are other evaluator types.
		\end{itemize}
	\end{frame}
	
	% ---------------------------------------------------------------
	\begin{frame}{Question 437}
		What does MLflow's model validation feature do?
		\begin{enumerate}[label=\Alph*.]
			\item Compares a candidate model to a baseline
			\item Trains a new model
			\item Registers a model
			\item Serves a model
		\end{enumerate}
	\end{frame}
	\begin{frame}{Answer 437}
		\textbf{Correct: A} --- validation compares against a baseline.\\
		\vspace{0.5em}
		\textbf{Wrong:}
		\begin{itemize}
			\item B, C, D are unrelated.
		\end{itemize}
	\end{frame}
	
	% ---------------------------------------------------------------
	\begin{frame}{Question 438}
		What is a baseline model in validation?
		\begin{enumerate}[label=\Alph*.]
			\item A reference model used for comparison
			\item The newest model
			\item The training data
			\item The cluster
		\end{enumerate}
	\end{frame}
	\begin{frame}{Answer 438}
		\textbf{Correct: A} --- a baseline is a reference model.\\
		\vspace{0.5em}
		\textbf{Wrong:}
		\begin{itemize}
			\item B, C, D are not baselines.
		\end{itemize}
	\end{frame}
	
	% ---------------------------------------------------------------
	\begin{frame}{Question 439}
		What is the purpose of monitoring model performance in production?
		\begin{enumerate}[label=\Alph*.]
			\item Detecting degradation and drift
			\item Training new models
			\item Registering models
			\item Deleting models
		\end{enumerate}
	\end{frame}
	\begin{frame}{Answer 439}
		\textbf{Correct: A} --- monitoring detects degradation.\\
		\vspace{0.5em}
		\textbf{Wrong:}
		\begin{itemize}
			\item B, C, D are unrelated.
		\end{itemize}
	\end{frame}
	
	% ---------------------------------------------------------------
	\begin{frame}{Question 440}
		What metric is typically monitored for model drift?
		\begin{enumerate}[label=\Alph*.]
			\item Population stability index (PSI) or distribution distance
			\item Training loss
			\item Cluster CPU
			\item Storage size
		\end{enumerate}
	\end{frame}
	\begin{frame}{Answer 440}
		\textbf{Correct: A} --- PSI or distribution distances detect drift.\\
		\vspace{0.5em}
		\textbf{Wrong:}
		\begin{itemize}
			\item B, C, D are not drift metrics.
		\end{itemize}
	\end{frame}
	
	% ---------------------------------------------------------------
	\begin{frame}{Question 441}
		What is a data quality expectation?
		\begin{enumerate}[label=\Alph*.]
			\item A rule that data should satisfy
			\item A model
			\item A feature table
			\item A notebook
		\end{enumerate}
	\end{frame}
	\begin{frame}{Answer 441}
		\textbf{Correct: A} --- expectations are data rules.\\
		\vspace{0.5em}
		\textbf{Wrong:}
		\begin{itemize}
			\item B, C, D are unrelated.
		\end{itemize}
	\end{frame}
	
	% ---------------------------------------------------------------
	\begin{frame}{Question 442}
		What is a model explainability method?
		\begin{enumerate}[label=\Alph*.]
			\item SHAP
			\item PCA
			\item KMeans
			\item ALS
		\end{enumerate}
	\end{frame}
	\begin{frame}{Answer 442}
		\textbf{Correct: A} --- SHAP explains model predictions.\\
		\vspace{0.5em}
		\textbf{Wrong:}
		\begin{itemize}
			\item B, C, D are other algorithms.
		\end{itemize}
	\end{frame}
	
	% ---------------------------------------------------------------
	\begin{frame}{Question 443}
		What is a permutation feature importance?
		\begin{enumerate}[label=\Alph*.]
			\item A method that measures accuracy drop when a feature is shuffled
			\item A model
			\item A feature table
			\item A metric
		\end{enumerate}
	\end{frame}
	\begin{frame}{Answer 443}
		\textbf{Correct: A} --- permutation importance measures feature impact.\\
		\vspace{0.5em}
		\textbf{Wrong:}
		\begin{itemize}
			\item B, C, D are unrelated.
		\end{itemize}
	\end{frame}
	
	% ---------------------------------------------------------------
	\begin{frame}{Question 444}
		What is partial dependence?
		\begin{enumerate}[label=\Alph*.]
			\item The marginal effect of a feature on predictions
			\item A model type
			\item A feature table
			\item A metric
		\end{enumerate}
	\end{frame}
	\begin{frame}{Answer 444}
		\textbf{Correct: A} --- partial dependence shows marginal effects.\\
		\vspace{0.5em}
		\textbf{Wrong:}
		\begin{itemize}
			\item B, C, D are unrelated.
		\end{itemize}
	\end{frame}
	
	% ---------------------------------------------------------------
	\begin{frame}{Question 445}
		What is distributed training?
		\begin{enumerate}[label=\Alph*.]
			\item Training across multiple machines or devices
			\item Training a single model on one node
			\item Serving a model
			\item Registering a model
		\end{enumerate}
	\end{frame}
	\begin{frame}{Answer 445}
		\textbf{Correct: A} --- distributed training uses multiple resources.\\
		\vspace{0.5em}
		\textbf{Wrong:}
		\begin{itemize}
			\item B, C, D are other tasks.
		\end{itemize}
	\end{frame}
	
	% ---------------------------------------------------------------
	\begin{frame}{Question 446}
		Which API enables distributed training in Spark ML?
		\begin{enumerate}[label=\Alph*.]
			\item \texttt{Spark ML estimators on Spark DataFrames}
			\item \texttt{pandas UDFs only}
			\item \texttt{Notebook widgets}
			\item \texttt{SQL warehouses}
		\end{enumerate}
	\end{frame}
	\begin{frame}{Answer 446}
		\textbf{Correct: A} --- Spark ML trains across a cluster.\\
		\vspace{0.5em}
		\textbf{Wrong:}
		\begin{itemize}
			\item B, C, D are not distributed training APIs.
		\end{itemize}
	\end{frame}
	
	% ---------------------------------------------------------------
	\begin{frame}{Question 447}
		What is a HorovodRunner used for in Databricks?
		\begin{enumerate}[label=\Alph*.]
			\item Distributed deep learning training
			\item SQL queries
			\item Model serving
			\item Data storage
		\end{enumerate}
	\end{frame}
	\begin{frame}{Answer 447}
		\textbf{Correct: A} --- HorovodRunner enables distributed deep learning.\\
		\vspace{0.5em}
		\textbf{Wrong:}
		\begin{itemize}
			\item B, C, D are unrelated.
		\end{itemize}
	\end{frame}
	
	% ---------------------------------------------------------------
	\begin{frame}{Question 448}
		What is the main benefit of distributed training?
		\begin{enumerate}[label=\Alph*.]
			\item Reduced training time for large data or models
			\item Lower accuracy
			\item Fewer features
			\item Smaller storage
		\end{enumerate}
	\end{frame}
	\begin{frame}{Answer 448}
		\textbf{Correct: A} --- distribution reduces training time.\\
		\vspace{0.5em}
		\textbf{Wrong:}
		\begin{itemize}
			\item B, C, D are not benefits.
		\end{itemize}
	\end{frame}
	
	% ---------------------------------------------------------------
	\begin{frame}{Question 449}
		What is a common challenge in distributed training?
		\begin{enumerate}[label=\Alph*.]
			\item Communication overhead between workers
			\item Too little data
			\item Too few models
			\item Too few metrics
		\end{enumerate}
	\end{frame}
	\begin{frame}{Answer 449}
		\textbf{Correct: A} --- communication overhead is a key challenge.\\
		\vspace{0.5em}
		\textbf{Wrong:}
		\begin{itemize}
			\item B, C, D are not the challenge.
		\end{itemize}
	\end{frame}
	
	% ---------------------------------------------------------------
	\begin{frame}{Question 450}
		What is a data parallel training strategy?
		\begin{enumerate}[label=\Alph*.]
			\item Splitting data across workers
			\item Splitting the model across workers
			\item Splitting metrics
			\item Splitting clusters
		\end{enumerate}
	\end{frame}
	\begin{frame}{Answer 450}
		\textbf{Correct: A} --- data parallel splits data.\\
		\vspace{0.5em}
		\textbf{Wrong:}
		\begin{itemize}
			\item B is model parallel.
			\item C, D are not strategies.
		\end{itemize}
	\end{frame}
	
	% ---------------------------------------------------------------
	\begin{frame}{Question 451}
		What is a model parallel training strategy?
		\begin{enumerate}[label=\Alph*.]
			\item Splitting the model across workers
			\item Splitting data across workers
			\item Splitting metrics
			\item Splitting clusters
		\end{enumerate}
	\end{frame}
	\begin{frame}{Answer 451}
		\textbf{Correct: A} --- model parallel splits the model.\\
		\vspace{0.5em}
		\textbf{Wrong:}
		\begin{itemize}
			\item B is data parallel.
			\item C, D are not strategies.
		\end{itemize}
	\end{frame}
	
	% ---------------------------------------------------------------
	\begin{frame}{Question 452}
		What is a checkpoint in distributed training?
		\begin{enumerate}[label=\Alph*.]
			\item A saved model state used for recovery or resumption
			\item A metric
			\item A cluster
			\item A feature table
		\end{enumerate}
	\end{frame}
	\begin{frame}{Answer 452}
		\textbf{Correct: A} --- checkpoints save model state.\\
		\vspace{0.5em}
		\textbf{Wrong:}
		\begin{itemize}
			\item B, C, D are unrelated.
		\end{itemize}
	\end{frame}
	
	% ---------------------------------------------------------------
	\begin{frame}{Question 453}
		Why are checkpoints useful for long training jobs?
		\begin{enumerate}[label=\Alph*.]
			\item They allow resuming after failure
			\item They reduce features
			\item They register models
			\item They serve models
		\end{enumerate}
	\end{frame}
	\begin{frame}{Answer 453}
		\textbf{Correct: A} --- checkpoints enable resumption.\\
		\vspace{0.5em}
		\textbf{Wrong:}
		\begin{itemize}
			\item B, C, D are unrelated.
		\end{itemize}
	\end{frame}
	
	% ---------------------------------------------------------------
	\begin{frame}{Question 454}
		What is early stopping in training?
		\begin{enumerate}[label=\Alph*.]
			\item Ending training when validation performance stops improving
			\item Ending training after one epoch
			\item Ending training at a fixed iteration
			\item Ending training when training error is zero
		\end{enumerate}
	\end{frame}
	\begin{frame}{Answer 454}
		\textbf{Correct: A} --- early stopping uses validation performance.\\
		\vspace{0.5em}
		\textbf{Wrong:}
		\begin{itemize}
			\item B, C, D are arbitrary stopping rules.
		\end{itemize}
	\end{frame}
	
	% ---------------------------------------------------------------
	\begin{frame}{Question 455}
		What is a learning rate in optimization?
		\begin{enumerate}[label=\Alph*.]
			\item The step size for parameter updates
			\item The number of trees
			\item The number of features
			\item The number of rows
		\end{enumerate}
	\end{frame}
	\begin{frame}{Answer 455}
		\textbf{Correct: A} --- learning rate controls step size.\\
		\vspace{0.5em}
		\textbf{Wrong:}
		\begin{itemize}
			\item B, C, D are not learning rates.
		\end{itemize}
	\end{frame}
	
	% ---------------------------------------------------------------
	\begin{frame}{Question 456}
		What happens if the learning rate is too high?
		\begin{enumerate}[label=\Alph*.]
			\item Training may diverge or overshoot
			\item Training becomes stable
			\item Accuracy always improves
			\item Storage increases
		\end{enumerate}
	\end{frame}
	\begin{frame}{Answer 456}
		\textbf{Correct: A} --- a high learning rate can diverge.\\
		\vspace{0.5em}
		\textbf{Wrong:}
		\begin{itemize}
			\item B, C, D are false.
		\end{itemize}
	\end{frame}
	
	% ---------------------------------------------------------------
	\begin{frame}{Question 457}
		What happens if the learning rate is too low?
		\begin{enumerate}[label=\Alph*.]
			\item Training becomes very slow
			\item Training diverges
			\item Accuracy always improves
			\item Storage increases
		\end{enumerate}
	\end{frame}
	\begin{frame}{Answer 457}
		\textbf{Correct: A} --- low learning rates train slowly.\\
		\vspace{0.5em}
		\textbf{Wrong:}
		\begin{itemize}
			\item B, C, D are false.
		\end{itemize}
	\end{frame}
	
	% ---------------------------------------------------------------
	\begin{frame}{Question 458}
		What is a mini-batch in training?
		\begin{enumerate}[label=\Alph*.]
			\item A small subset of data used for one update
			\item The full dataset
			\item A cluster
			\item A metric
		\end{enumerate}
	\end{frame}
	\begin{frame}{Answer 458}
		\textbf{Correct: A} --- mini-batches update parameters incrementally.\\
		\vspace{0.5em}
		\textbf{Wrong:}
		\begin{itemize}
			\item B, C, D are unrelated.
		\end{itemize}
	\end{frame}
	
	% ---------------------------------------------------------------
	\begin{frame}{Question 459}
		What is an epoch in training?
		\begin{enumerate}[label=\Alph*.]
			\item One full pass over the training data
			\item One update step
			\item One metric
			\item One cluster
		\end{enumerate}
	\end{frame}
	\begin{frame}{Answer 459}
		\textbf{Correct: A} --- an epoch is a full data pass.\\
		\vspace{0.5em}
		\textbf{Wrong:}
		\begin{itemize}
			\item B, C, D are unrelated.
		\end{itemize}
	\end{frame}
	
	% ---------------------------------------------------------------
	\begin{frame}{Question 460}
		What is a batch size in training?
		\begin{enumerate}[label=\Alph*.]
			\item The number of samples per update
			\item The number of epochs
			\item The number of trees
			\item The number of features
		\end{enumerate}
	\end{frame}
	\begin{frame}{Answer 460}
		\textbf{Correct: A} --- batch size sets samples per update.\\
		\vspace{0.5em}
		\textbf{Wrong:}
		\begin{itemize}
			\item B, C, D are other settings.
		\end{itemize}
	\end{frame}
	
	% ---------------------------------------------------------------
	\begin{frame}{Question 461}
		What is a confusion matrix useful for beyond accuracy?
		\begin{enumerate}[label=\Alph*.]
			\item It shows false positives and false negatives separately
			\item It trains a model
			\item It stores data
			\item It serves a model
		\end{enumerate}
	\end{frame}
	\begin{frame}{Answer 461}
		\textbf{Correct: A} --- it separates FP and FN counts.\\
		\vspace{0.5em}
		\textbf{Wrong:}
		\begin{itemize}
			\item B, C, D are unrelated.
		\end{itemize}
	\end{frame}
	
	% ---------------------------------------------------------------
	\begin{frame}{Question 462}
		What is a ROC curve?
		\begin{enumerate}[label=\Alph*.]
			\item A plot of true positive rate versus false positive rate
			\item A plot of precision versus recall
			\item A plot of error versus training size
			\item A plot of feature importance
		\end{enumerate}
	\end{frame}
	\begin{frame}{Answer 462}
		\textbf{Correct: A} --- ROC plots TPR vs.\ FPR.\\
		\vspace{0.5em}
		\textbf{Wrong:}
		\begin{itemize}
			\item B is a PR curve.
			\item C is a learning curve.
			\item D is feature importance.
		\end{itemize}
	\end{frame}
	
	% ---------------------------------------------------------------
	\begin{frame}{Question 463}
		What is a PR curve?
		\begin{enumerate}[label=\Alph*.]
			\item A plot of precision versus recall
			\item A plot of TPR versus FPR
			\item A plot of error versus training size
			\item A plot of feature importance
		\end{enumerate}
	\end{frame}
	\begin{frame}{Answer 463}
		\textbf{Correct: A} --- PR curves plot precision vs.\ recall.\\
		\vspace{0.5em}
		\textbf{Wrong:}
		\begin{itemize}
			\item B is ROC.
			\item C is a learning curve.
			\item D is feature importance.
		\end{itemize}
	\end{frame}
	
	% ---------------------------------------------------------------
	\begin{frame}{Question 464}
		What is log loss?
		\begin{enumerate}[label=\Alph*.]
			\item A classification metric that penalizes confident wrong predictions
			\item A regression metric
			\item A clustering metric
			\item A drift metric
		\end{enumerate}
	\end{frame}
	\begin{frame}{Answer 464}
		\textbf{Correct: A} --- log loss penalizes confident errors.\\
		\vspace{0.5em}
		\textbf{Wrong:}
		\begin{itemize}
			\item B, C, D are other metrics.
		\end{itemize}
	\end{frame}
	
	% ---------------------------------------------------------------
	\begin{frame}{Question 465}
		What is cross-entropy?
		\begin{enumerate}[label=\Alph*.]
			\item A loss function for classification
			\item A clustering metric
			\item A drift metric
			\item A regression metric
		\end{enumerate}
	\end{frame}
	\begin{frame}{Answer 465}
		\textbf{Correct: A} --- cross-entropy is a classification loss.\\
		\vspace{0.5em}
		\textbf{Wrong:}
		\begin{itemize}
			\item B, C, D are other metrics.
		\end{itemize}
	\end{frame}
	
	% ---------------------------------------------------------------
	\begin{frame}{Question 466}
		What is a target variable in supervised learning?
		\begin{enumerate}[label=\Alph*.]
			\item The value being predicted
			\item The input features
			\item The cluster
			\item The model
		\end{enumerate}
	\end{frame}
	\begin{frame}{Answer 466}
		\textbf{Correct: A} --- the target is predicted.\\
		\vspace{0.5em}
		\textbf{Wrong:}
		\begin{itemize}
			\item B, C, D are unrelated.
		\end{itemize}
	\end{frame}
	
	% ---------------------------------------------------------------
	\begin{frame}{Question 467}
		What is a feature in supervised learning?
		\begin{enumerate}[label=\Alph*.]
			\item An input variable used to predict the target
			\item The predicted value
			\item The model
			\item The cluster
		\end{enumerate}
	\end{frame}
	\begin{frame}{Answer 467}
		\textbf{Correct: A} --- features are input variables.\\
		\vspace{0.5em}
		\textbf{Wrong:}
		\begin{itemize}
			\item B, C, D are unrelated.
		\end{itemize}
	\end{frame}
	
	% ---------------------------------------------------------------
	\begin{frame}{Question 468}
		What is a label in supervised learning?
		\begin{enumerate}[label=\Alph*.]
			\item The known target value in training data
			\item The input features
			\item The model
			\item The cluster
		\end{enumerate}
	\end{frame}
	\begin{frame}{Answer 468}
		\textbf{Correct: A} --- labels are known target values.\\
		\vspace{0.5em}
		\textbf{Wrong:}
		\begin{itemize}
			\item B, C, D are unrelated.
		\end{itemize}
	\end{frame}
	
	% ---------------------------------------------------------------
	\begin{frame}{Question 469}
		What is semi-supervised learning?
		\begin{enumerate}[label=\Alph*.]
			\item Learning from a mix of labeled and unlabeled data
			\item Learning from only labeled data
			\item Learning from only unlabeled data
			\item Learning without a model
		\end{enumerate}
	\end{frame}
	\begin{frame}{Answer 469}
		\textbf{Correct: A} --- semi-supervised uses both labeled and unlabeled data.\\
		\vspace{0.5em}
		\textbf{Wrong:}
		\begin{itemize}
			\item B is supervised.
			\item C is unsupervised.
			\item D is false.
		\end{itemize}
	\end{frame}
	
	% ---------------------------------------------------------------
	\begin{frame}{Question 470}
		What is self-supervised learning?
		\begin{enumerate}[label=\Alph*.]
			\item Learning from data-generated labels
			\item Learning from manual labels
			\item Learning from clusters
			\item Learning without data
		\end{enumerate}
	\end{frame}
	\begin{frame}{Answer 470}
		\textbf{Correct: A} --- self-supervised uses generated labels.\\
		\vspace{0.5em}
		\textbf{Wrong:}
		\begin{itemize}
			\item B is supervised.
			\item C is unsupervised.
			\item D is false.
		\end{itemize}
	\end{frame}
	
	% ---------------------------------------------------------------
	\begin{frame}{Question 471}
		What is reinforcement learning?
		\begin{enumerate}[label=\Alph*.]
			\item Learning from rewards and penalties
			\item Learning from labeled examples
			\item Learning from clusters
			\item Learning from no data
		\end{enumerate}
	\end{frame}
	\begin{frame}{Answer 471}
		\textbf{Correct: A} --- RL uses rewards and penalties.\\
		\vspace{0.5em}
		\textbf{Wrong:}
		\begin{itemize}
			\item B is supervised.
			\item C is unsupervised.
			\item D is false.
		\end{itemize}
	\end{frame}
	
	% ---------------------------------------------------------------
	\begin{frame}{Question 472}
		What is a reward signal in reinforcement learning?
		\begin{enumerate}[label=\Alph*.]
			\item Feedback guiding the agent's behavior
			\item A feature
			\item A label
			\item A cluster
		\end{enumerate}
	\end{frame}
	\begin{frame}{Answer 472}
		\textbf{Correct: A} --- rewards guide learning.\\
		\vspace{0.5em}
		\textbf{Wrong:}
		\begin{itemize}
			\item B, C, D are unrelated.
		\end{itemize}
	\end{frame}
	
	% ---------------------------------------------------------------
	\begin{frame}{Question 473}
		What is a policy in reinforcement learning?
		\begin{enumerate}[label=\Alph*.]
			\item A mapping from states to actions
			\item A feature
			\item A label
			\item A cluster
		\end{enumerate}
	\end{frame}
	\begin{frame}{Answer 473}
		\textbf{Correct: A} --- policies map states to actions.\\
		\vspace{0.5em}
		\textbf{Wrong:}
		\begin{itemize}
			\item B, C, D are unrelated.
		\end{itemize}
	\end{frame}
	
	% ---------------------------------------------------------------
	\begin{frame}{Question 474}
		What is an agent in reinforcement learning?
		\begin{enumerate}[label=\Alph*.]
			\item The learner that takes actions
			\item The environment
			\item The reward
			\item The state
		\end{enumerate}
	\end{frame}
	\begin{frame}{Answer 474}
		\textbf{Correct: A} --- the agent takes actions.\\
		\vspace{0.5em}
		\textbf{Wrong:}
		\begin{itemize}
			\item B, C, D are other RL components.
		\end{itemize}
	\end{frame}
	
	% ---------------------------------------------------------------
	\begin{frame}{Question 475}
		What is an environment in reinforcement learning?
		\begin{enumerate}[label=\Alph*.]
			\item The world the agent interacts with
			\item The agent
			\item The reward
			\item The policy
		\end{enumerate}
	\end{frame}
	\begin{frame}{Answer 475}
		\textbf{Correct: A} --- the environment is the world.\\
		\vspace{0.5em}
		\textbf{Wrong:}
		\begin{itemize}
			\item B, C, D are other RL components.
		\end{itemize}
	\end{frame}
	
	% ---------------------------------------------------------------
	\begin{frame}{Question 476}
		What is an observation in reinforcement learning?
		\begin{enumerate}[label=\Alph*.]
			\item Information the agent perceives about the environment
			\item The reward
			\item The policy
			\item The model
		\end{enumerate}
	\end{frame}
	\begin{frame}{Answer 476}
		\textbf{Correct: A} --- observations are perceived information.\\
		\vspace{0.5em}
		\textbf{Wrong:}
		\begin{itemize}
			\item B, C, D are other components.
		\end{itemize}
	\end{frame}
	
	% ---------------------------------------------------------------
	\begin{frame}{Question 477}
		What is an action in reinforcement learning?
		\begin{enumerate}[label=\Alph*.]
			\item A choice the agent makes
			\item A reward
			\item A policy
			\item A model
		\end{enumerate}
	\end{frame}
	\begin{frame}{Answer 477}
		\textbf{Correct: A} --- actions are agent choices.\\
		\vspace{0.5em}
		\textbf{Wrong:}
		\begin{itemize}
			\item B, C, D are other components.
		\end{itemize}
	\end{frame}
	
	% ---------------------------------------------------------------
	\begin{frame}{Question 478}
		What is exploration in reinforcement learning?
		\begin{enumerate}[label=\Alph*.]
			\item Trying new actions to learn
			\item Repeating known actions
			\item Training a model
			\item Serving a model
		\end{enumerate}
	\end{frame}
	\begin{frame}{Answer 478}
		\textbf{Correct: A} --- exploration tries new actions.\\
		\vspace{0.5em}
		\textbf{Wrong:}
		\begin{itemize}
			\item B is exploitation.
			\item C, D are unrelated.
		\end{itemize}
	\end{frame}
	
	% ---------------------------------------------------------------
	\begin{frame}{Question 479}
		What is exploitation in reinforcement learning?
		\begin{enumerate}[label=\Alph*.]
			\item Choosing the best-known action
			\item Trying new actions
			\item Training a model
			\item Serving a model
		\end{enumerate}
	\end{frame}
	\begin{frame}{Answer 479}
		\textbf{Correct: A} --- exploitation uses known best actions.\\
		\vspace{0.5em}
		\textbf{Wrong:}
		\begin{itemize}
			\item B is exploration.
			\item C, D are unrelated.
		\end{itemize}
	\end{frame}
	
	% ---------------------------------------------------------------
	\begin{frame}{Question 480}
		What is the exploration-exploitation tradeoff?
		\begin{enumerate}[label=\Alph*.]
			\item Balancing trying new actions versus using known good actions
			\item Balancing training and serving
			\item Balancing data and models
			\item Balancing clusters
		\end{enumerate}
	\end{frame}
	\begin{frame}{Answer 480}
		\textbf{Correct: A} --- the tradeoff balances exploration and exploitation.\\
		\vspace{0.5em}
		\textbf{Wrong:}
		\begin{itemize}
			\item B, C, D are unrelated.
		\end{itemize}
	\end{frame}
	
	% ---------------------------------------------------------------
	\begin{frame}{Question 481}
		What is a Markov decision process?
		\begin{enumerate}[label=\Alph*.]
			\item A mathematical framework for RL with states, actions, and rewards
			\item A model type
			\item A feature table
			\item A cluster
		\end{enumerate}
	\end{frame}
	\begin{frame}{Answer 481}
		\textbf{Correct: A} --- an MDP defines RL states/actions/rewards.\\
		\vspace{0.5em}
		\textbf{Wrong:}
		\begin{itemize}
			\item B, C, D are unrelated.
		\end{itemize}
	\end{frame}
	
	% ---------------------------------------------------------------
	\begin{frame}{Question 482}
		What is a state in reinforcement learning?
		\begin{enumerate}[label=\Alph*.]
			\item A representation of the environment at a point in time
			\item An action
			\item A reward
			\item A policy
		\end{enumerate}
	\end{frame}
	\begin{frame}{Answer 482}
		\textbf{Correct: A} --- states represent the environment.\\
		\vspace{0.5em}
		\textbf{Wrong:}
		\begin{itemize}
			\item B, C, D are other components.
		\end{itemize}
	\end{frame}
	
	% ---------------------------------------------------------------
	\begin{frame}{Question 483}
		What is a value function in reinforcement learning?
		\begin{enumerate}[label=\Alph*.]
			\item An estimate of expected future reward
			\item A feature
			\item A label
			\item A cluster
		\end{enumerate}
	\end{frame}
	\begin{frame}{Answer 483}
		\textbf{Correct: A} --- value functions estimate future reward.\\
		\vspace{0.5em}
		\textbf{Wrong:}
		\begin{itemize}
			\item B, C, D are unrelated.
		\end{itemize}
	\end{frame}
	
	% ---------------------------------------------------------------
	\begin{frame}{Question 484}
		What is Q-learning?
		\begin{enumerate}[label=\Alph*.]
			\item A reinforcement learning algorithm that learns action values
			\item A clustering algorithm
			\item A regression algorithm
			\item A dimensionality reduction algorithm
		\end{enumerate}
	\end{frame}
	\begin{frame}{Answer 484}
		\textbf{Correct: A} --- Q-learning learns action values.\\
		\vspace{0.5em}
		\textbf{Wrong:}
		\begin{itemize}
			\item B, C, D are other algorithms.
		\end{itemize}
	\end{frame}
	
	% ---------------------------------------------------------------
	\begin{frame}{Question 485}
		What is a policy gradient method?
		\begin{enumerate}[label=\Alph*.]
			\item A method that optimizes the policy directly
			\item A method that clusters data
			\item A method that reduces features
			\item A method that serves models
		\end{enumerate}
	\end{frame}
	\begin{frame}{Answer 485}
		\textbf{Correct: A} --- policy gradients optimize policies directly.\\
		\vspace{0.5em}
		\textbf{Wrong:}
		\begin{itemize}
			\item B, C, D are unrelated.
		\end{itemize}
	\end{frame}
	
	% ---------------------------------------------------------------
	\begin{frame}{Question 486}
		What is a bandit problem?
		\begin{enumerate}[label=\Alph*.]
			\item A simplified RL problem with one state and multiple actions
			\item A classification problem
			\item A regression problem
			\item A clustering problem
		\end{enumerate}
	\end{frame}
	\begin{frame}{Answer 486}
		\textbf{Correct: A} --- bandits are single-state RL problems.\\
		\vspace{0.5em}
		\textbf{Wrong:}
		\begin{itemize}
			\item B, C, D are other problem types.
		\end{itemize}
	\end{frame}
	
	% ---------------------------------------------------------------
	\begin{frame}{Question 487}
		What is A/B testing in the context of bandits?
		\begin{enumerate}[label=\Alph*.]
			\item Comparing two actions using observed rewards
			\item Training two models
			\item Deleting two models
			\item Registering two models
		\end{enumerate}
	\end{frame}
	\begin{frame}{Answer 487}
		\textbf{Correct: A} --- A/B tests compare actions.\\
		\vspace{0.5em}
		\textbf{Wrong:}
		\begin{itemize}
			\item B, C, D are unrelated.
		\end{itemize}
	\end{frame}
	
	% ---------------------------------------------------------------
	\begin{frame}{Question 488}
		What is a multi-armed bandit?
		\begin{enumerate}[label=\Alph*.]
			\item A problem of choosing among multiple actions with unknown rewards
			\item A clustering problem
			\item A regression problem
			\item A dimensionality reduction problem
		\end{enumerate}
	\end{frame}
	\begin{frame}{Answer 488}
		\textbf{Correct: A} --- multi-armed bandits involve multiple actions.\\
		\vspace{0.5em}
		\textbf{Wrong:}
		\begin{itemize}
			\item B, C, D are other problem types.
		\end{itemize}
	\end{frame}
	
	% ---------------------------------------------------------------
	\begin{frame}{Question 489}
		What is Thompson sampling?
		\begin{enumerate}[label=\Alph*.]
			\item A bandit algorithm using posterior distributions
			\item A clustering algorithm
			\item A regression algorithm
			\item A serving method
		\end{enumerate}
	\end{frame}
	\begin{frame}{Answer 489}
		\textbf{Correct: A} --- Thompson sampling uses posteriors.\\
		\vspace{0.5em}
		\textbf{Wrong:}
		\begin{itemize}
			\item B, C, D are unrelated.
		\end{itemize}
	\end{frame}
	
	% ---------------------------------------------------------------
	\begin{frame}{Question 490}
		What is epsilon-greedy?
		\begin{enumerate}[label=\Alph*.]
			\item A bandit strategy that explores with probability epsilon
			\item A clustering method
			\item A regression method
			\item A serving method
		\end{enumerate}
	\end{frame}
	\begin{frame}{Answer 490}
		\textbf{Correct: A} --- epsilon-greedy balances exploration/exploitation.\\
		\vspace{0.5em}
		\textbf{Wrong:}
		\begin{itemize}
			\item B, C, D are unrelated.
		\end{itemize}
	\end{frame}
	
	% ---------------------------------------------------------------
	\begin{frame}{Question 491}
		What is a feature store?
		\begin{enumerate}[label=\Alph*.]
			\item A centralized repository for features
			\item A model
			\item A cluster
			\item A notebook
		\end{enumerate}
	\end{frame}
	\begin{frame}{Answer 491}
		\textbf{Correct: A} --- a feature store centralizes features.\\
		\vspace{0.5em}
		\textbf{Wrong:}
		\begin{itemize}
			\item B, C, D are unrelated.
		\end{itemize}
	\end{frame}
	
	% ---------------------------------------------------------------
	\begin{frame}{Question 492}
		What is a feature view?
		\begin{enumerate}[label=\Alph*.]
			\item A logical view of features for a use case
			\item A model
			\item A cluster
			\item A table
		\end{enumerate}
	\end{frame}
	\begin{frame}{Answer 492}
		\textbf{Correct: A} --- feature views organize features.\\
		\vspace{0.5em}
		\textbf{Wrong:}
		\begin{itemize}
			\item B, C, D are unrelated.
		\end{itemize}
	\end{frame}
	
	% ---------------------------------------------------------------
	\begin{frame}{Question 493}
		What is a feature group?
		\begin{enumerate}[label=\Alph*.]
			\item A set of related features
			\item A model
			\item A cluster
			\item A notebook
		\end{enumerate}
	\end{frame}
	\begin{frame}{Answer 493}
		\textbf{Correct: A} --- feature groups are related features.\\
		\vspace{0.5em}
		\textbf{Wrong:}
		\begin{itemize}
			\item B, C, D are unrelated.
		\end{itemize}
	\end{frame}
	
	% ---------------------------------------------------------------
	\begin{frame}{Question 494}
		What is a feature pipeline?
		\begin{enumerate}[label=\Alph*.]
			\item A workflow that computes and updates features
			\item A model
			\item A cluster
			\item A notebook
		\end{enumerate}
	\end{frame}
	\begin{frame}{Answer 494}
		\textbf{Correct: A} --- feature pipelines compute features.\\
		\vspace{0.5em}
		\textbf{Wrong:}
		\begin{itemize}
			\item B, C, D are unrelated.
		\end{itemize}
	\end{frame}
	
	% ---------------------------------------------------------------
	\begin{frame}{Question 495}
		What is a feature transformation?
		\begin{enumerate}[label=\Alph*.]
			\item An operation that creates or modifies features
			\item A model
			\item A cluster
			\item A notebook
		\end{enumerate}
	\end{frame}
	\begin{frame}{Answer 495}
		\textbf{Correct: A} --- transformations modify features.\\
		\vspace{0.5em}
		\textbf{Wrong:}
		\begin{itemize}
			\item B, C, D are unrelated.
		\end{itemize}
	\end{frame}
	
	% ---------------------------------------------------------------
	\begin{frame}{Question 496}
		What is a feature backfill?
		\begin{enumerate}[label=\Alph*.]
			\item Computing historical feature values
			\item Training a model
			\item Serving a model
			\item Deleting features
		\end{enumerate}
	\end{frame}
	\begin{frame}{Answer 496}
		\textbf{Correct: A} --- backfills compute historical values.\\
		\vspace{0.5em}
		\textbf{Wrong:}
		\begin{itemize}
			\item B, C, D are unrelated.
		\end{itemize}
	\end{frame}
	
	% ---------------------------------------------------------------
	\begin{frame}{Question 497}
		Why is a feature backfill needed after a feature definition changes?
		\begin{enumerate}[label=\Alph*.]
			\item Historical feature values must be recomputed consistently
			\item The model retrains automatically
			\item The cluster resizes
			\item The notebook is deleted
		\end{enumerate}
	\end{frame}
	\begin{frame}{Answer 497}
		\textbf{Correct: A} --- backfills keep history consistent.\\
		\vspace{0.5em}
		\textbf{Wrong:}
		\begin{itemize}
			\item B, C, D are false.
		\end{itemize}
	\end{frame}
	
	% ---------------------------------------------------------------
	\begin{frame}{Question 498}
		What is feature versioning?
		\begin{enumerate}[label=\Alph*.]
			\item Tracking changes to feature definitions over time
			\item Training a model
			\item Serving a model
			\item Registering a model
		\end{enumerate}
	\end{frame}
	\begin{frame}{Answer 498}
		\textbf{Correct: A} --- feature versioning tracks definition changes.\\
		\vspace{0.5em}
		\textbf{Wrong:}
		\begin{itemize}
			\item B, C, D are unrelated.
		\end{itemize}
	\end{frame}
	
	% ---------------------------------------------------------------
	\begin{frame}{Question 499}
		What is a feature store's primary benefit for collaboration?
		\begin{enumerate}[label=\Alph*.]
			\item Shared, governed feature definitions across teams
			\item Faster training
			\item Smaller models
			\item Fewer clusters
		\end{enumerate}
	\end{frame}
	\begin{frame}{Answer 499}
		\textbf{Correct: A} --- feature stores share governed features.\\
		\vspace{0.5em}
		\textbf{Wrong:}
		\begin{itemize}
			\item B, C, D are false.
		\end{itemize}
	\end{frame}
	
	% ---------------------------------------------------------------
	\begin{frame}{Question 500}
		What is the best practice for updating a feature table used in production?
		\begin{enumerate}[label=\Alph*.]
			\item Backfill, validate, and then update downstream models
			\item Delete the table
			\item Stop monitoring
			\item Skip validation
		\end{enumerate}
	\end{frame}
	\begin{frame}{Answer 500}
		\textbf{Correct: A} --- backfill and validate before rollout.\\
		\vspace{0.5em}
		\textbf{Wrong:}
		\begin{itemize}
			\item B causes downtime.
			\item C removes safeguards.
			\item D risks bad data.
		\end{itemize}
	\end{frame}
	
\end{document}